\section{未来实验验证}
\label{sec:future}

本文不包含任何湿实验数据；本节说明如何用最小实验设计检验平台输出的候选假设，并给出判定标准。

\textbf{方向一：第 18 位碱基替换（候选 C1）.} 选择若干携带目标位点的 sgRNA，仅替换第 18 位碱基（C$\leftrightarrow$A），保持其余 22\,nt 完全不变，在同一实验体系中测定编辑效率并做配对比较。该设计直接检验"模型优先排序的核苷酸差异是否伴随实测效率变化"，并可通过在多于一个细胞系中重复，检验平台提示的上下文依赖性（HCT116/HeLa 的 C$-$A 差约 $+0.09$，而 HEK293T 为 $-0.015$）。读出指标建议与本文数据一致（归一化编辑效率），并预先登记样本量、重复次数与分层比较方案。

\textbf{方向二：PAM 邻近候选序列模式破坏（候选 C2/C3）.} 在包含目标 4\,nt 核心的 sgRNA 中替换该核心，同时匹配 GC 含量与长度以控制混杂，比较携带与不携带核心序列的实测效率。由于富集证据目前仅见于单一细胞系，实验应优先在该细胞系中进行，并在其它细胞系作为探索性重复。

\textbf{方向三：独立数据集与真正的泛化验证.} 本文的 LOCO 分支已按无同源泄漏的方式重跑（448 次运行），结果显示跨细胞系泛化为 0 或负值；后续应在独立编辑数据集（不同细胞类型或不同编辑系统）上检验：(i) 平台流程的可迁移性；(ii) 位置层级序列信号（第 18/20 位）是否复现；(iii) 环境增量效应是否在更高分辨率的环境特征下变得可检测。

\textbf{方向四：分离"分布漂移"与"训练集缩减".} 本文 LOCO 中两种因素混杂（\S\ref{sec:limitations}）。建议的对照设计为：在保持训练样本量不变的前提下构造无同源泄漏的训练池（例如对留出系同源序列做等量重采样替代），从而分离二者的贡献。

\textbf{预期输出与判定标准.} 上述实验的预期输出为"候选假设被支持/不被支持"的判定与效应量估计；判定应基于预先登记的统计方案（分层比较与多重检验校正），并使用与本文一致的指标口径。无论方向如何，实验都将用于评估"AI 候选优先级 $\rightarrow$ 实验验证"闭环的实用价值。
